% Seção 9: confronto dos resultados com as afirmações da literatura (cada afirmação com a página da fonte)
Each item below quotes a statement of the literature and gives the corresponding result of this study.
\begin{enumerate}
\item \emph{Atmospheric pore air.} \citet[p.~811]{sheng2008} consider the atmospheric air pressure ``approximately
true for most in situ conditions''; \citet[p.~292]{fredlund2006} restricts it to soils ``exposed to atmospheric
pressure'' with a continuous air phase. In our analyses the maximum air pressure reaches 13~kPa after 3~h of
9~mm/h on the $f_{100}$ slope at $63.4^\circ$, before the first vertex ponds (6~h), and 36--38~kPa at the end of the
24-h rainfalls of 9--80~mm/h on the same slope, when 83 to all 96 surface vertices are ponded; in the $f_{10}$ soil it
reaches 40~kPa
with 140~mm/h. With active air the critical intensity of the reference slope is about half of that with passive air
(Section~\ref{sec:decomp}). Once the surface ponds, the soil below it is no longer exposed to the atmosphere, the
situation excluded by \citet{fredlund2006}.

\item \emph{Is the passive-air assumption unrealistic?} \citet[p.~2]{khoei2021} state that it leads ``to unrealistic
outcomes'' and cite \citet{callari2009} among the works that adopt it; \citet[p.~491]{callari2009} obtain ``good
agreement with available experimental data'' in the Liakopoulos test with it. In our analyses both statements hold for
different processes: in drainage (Liakopoulos) the passive-air solution reproduces the measurements and our
two-phase solution drains more slowly than the test (Section~\ref{sec:valid}); in infiltration with a ponded surface
the active air changes the critical intensity by a factor of about two. Neither statement is general.

\item \emph{Coupled or uncoupled.} \citet[p.~2]{pradhan2022} report that coupled analyses ``have been reported to
provide a better prediction of factor of safety''; \citet[p.~315]{fredlund2006} that ``the uncoupled simultaneous
solution yields satisfactory results for geotechnical engineering purposes''. In our analyses the coupling changes
the critical intensity by 5--11\% without dilatancy and by 27--62\% with an associated flow rule
($\delta=0.01$): the answer depends on the flow rule, which neither study discusses.

\item \emph{Flow rule of the coupled analyses.} The coupled analyses of slope collapse of
\citet[p.~6]{wang2018mpm} and \citet[p.~5]{liu2021} use dilation angles of $0.1^\circ$ and $0^\circ$. In our coupled
analyses the transient collapses without dilatancy ($I_r\ge26\mmh$) and does not collapse with $\psi\ge\phi'/2$
(Section~\ref{sec:dilat}). Whether a coupled analysis predicts a collapse therefore depends on a parameter that these
studies did not vary.

\item \emph{The intensity that matters relative to $k_s$.} \citet[p.~14]{liu2021}, citing \citet{rahardjo2007} among
others, place the threshold of $q/k_{sat}$ beyond which the FS no longer decreases at ``about or slightly larger
than 1''. The threshold intensities of \citet[p.~1539 and Fig.~7]{rahardjo2007} are 140, 51 and 11~mm/h for the
$f_{10}$, $f_{50}$ and $f_{100}$ soils, i.e.\ $0.39k_s$, $1.42k_s$ and $3k_s$, so the summary of \citet{liu2021} does
not match the values of the study it cites. Our critical intensities are about $0.4k_s$ in the $f_{10}$ soil in
both models and $1.3k_s$ (model iii) and 0.7--1.0$k_s$ (model i) in the $f_{50}$ soil: the ratio depends on the
soil and on the model.

\item \emph{The lower plateau.} The lower asymptotes of \citet[Fig.~7]{rahardjo2007} are 0.81 and 0.82 for the
$f_{10}$ and $f_{50}$ soils and 1.45 for the $f_{100}$ soil. \citet[p.~9]{li2024} find that the FS of a slope
saturated by a steady infiltration ``is identical to that without considering suction'', and
\citet[p.~17]{cheng2024} relate the constant minimum FS above a certain intensity to the degree of saturation, which
``keeps constant''. In our model (iii) the steady state with the surface ponded gives 0.84--0.85 for the three soils,
0.02--0.04 above the asymptotes of \citet{rahardjo2007} for $f_{10}$ and $f_{50}$; the $f_{100}$ slope does not reach
that state in 24~h, so its asymptote of 1.45 is not the FS of the saturated slope (Section~\ref{sec:criterio}).

\item \emph{Low-permeability slopes.} \citet[p.~1541]{rahardjo2007}: slopes with $k_s\le10^{-6}$~m/s ``will not fail
under short-duration rainfall ($T_r\le24$~h) even when the rainfall intensity is extremely high''. Our model (iii)
agrees (minimum FS 1.093 at $63.4^\circ$). With active air the minimum FS of the same slope is 1.013--1.015 (1.009 with
the finest mesh): failure with $\delta=0.02$, not with $\delta=0.01$ (Section~\ref{sec:ar}).

\item \emph{The minimum FS as the measure of safety.} \citet[p.~324]{tang2018} adopt the minimum FS during and after
the rainfall ``to reflect the most unstable state of the slope''. In our coupled models with an associated flow rule
the minimum FS did not fall below one in the cases computed (Section~\ref{sec:fs}), so beyond failure it does not distinguish the
rainfalls; models with different coupling can only be compared by the rainfall that brings the FS to one.
\end{enumerate}
